% ============================================================
%  第五章  社会科学与经济学中的因果性与结构模型
%  5.4—5.6 及章末
% ============================================================

\section{若干概念基础}

\subsection{结构参数究竟意味着什么？}

每个 SEM 的学生在生涯中的某个时刻，都会撞上下面这个悖论。若我们把方程
\[
y=\beta x+\varepsilon
\]
中的系数 $\beta$ 解释为 $X$ 每单位变化所引起的 $E(Y)$ 的变化，那么把方程改写为
\[
x=(y-\varepsilon)/\beta
\]
之后，我们就应当把 $1/\beta$ 解释为 $Y$ 每单位变化所引起的 $E(X)$ 的变化。但这既与直觉相悖，也与模型的预测相悖：若 $Y$ 不作为自变量出现在原来关于 $X$ 的结构方程中，则 $E(X)$ 随 $Y$ 的单位变化而发生的变化应当为零。

SEM 的教师通常经由两条逃逸路线之一来回避这一两难。一条路线否认 $\beta$ 具有任何因果解读，满足于纯统计的解释：$\beta$ 度量 $X$ 所解释的 $Y$ 的方差缩减（例如见 Muthen 1987）。另一条路线只允许那些满足“隔离”限制的系数具有因果解读（Bollen 1989; James et al. 1982）：解释变量必须与方程中的误差不相关。由于 $\varepsilon$ 不可能同时与 $X$ 和 $Y$ 都不相关（这种论证如此说），$\beta$ 与 $1/\beta$ 便不可能都具有因果含义，悖论也就消解了。

第一条路线是自洽的，但它损害了创立者们的本意——让 SEM 充当政策制定的辅助工具，也与多数 SEM 使用者的直觉相冲突。第二条路线在逻辑上易受攻击。众所周知，每一对二元正态变量 $X$ 与 $Y$ 都可以用两种等价的方式表达：
\[
y=\beta x+\varepsilon_1
\quad\text{与}\quad
x=\gamma y+\varepsilon_2,
\]
其中 $\mathrm{cov}(X,\varepsilon_1)=\mathrm{cov}(Y,\varepsilon_2)=0$，且 $\gamma=\rho_{XY}=\beta\sigma_X^2/\sigma_Y^2$。因此，若条件 $\mathrm{cov}(X,\varepsilon_1)=0$ 赋予 $\beta$ 以因果含义，那么 $\mathrm{cov}(Y,\varepsilon_2)=0$ 也应当赋予 $\gamma$ 以因果含义。但这也既与直觉相悖，也与 SEM 背后的意图相悖：若不存在从 $Y$ 到 $X$ 的因果路径，则 $E(X)$ 随 $Y$ 的单位变化而发生的变化应当为零，而不是 $\rho_{XY}$。

那么，结构系数究竟意味着什么？结构方程呢？误差项呢？因果效应的干预解释一旦与 $\dox{x}$ 记号相结合，便为这些问题提供了简单的答案。这些答案阐明了结构方程的操作含义，因而我希望它们能够终结一个围绕这些实体而生的争议与混乱的时代。

\subsubsection*{结构方程：操作定义}

\begin{definition}[结构方程]
方程 $y=\beta x+\varepsilon$ 被称为\emph{结构的}，若它应作如下解释：在一个理想实验中，我们控制 $X$ 取 $x$、并控制任一其他变量集 $Z$（不含 $X$ 或 $Y$）取 $z$，则 $Y$ 的取值 $y$ 由 $\beta x+\varepsilon$ 给出，其中 $\varepsilon$ 不是设定 $x$ 与 $z$ 的函数。
\end{definition}

这一定义是操作性的，因为一切量都是可观测的——尽管是在受控操纵的条件下。多数观测研究中无法实施操纵，这并不否定该定义的操作性；正如我们不能用肉眼观察细菌，并不否定细菌在显微镜下的可观测性。SEM 的挑战在于：从我们实际所能观测到的少量信息中，提取出关于我们希望观测之物的最大信息量。

注意，刚刚给出的操作解读并未对“当我们控制 $Y$ 时 $X$（或任何其他变量）会如何表现”作出任何断言。这种不对称性使结构方程中的等号不同于代数等号：在关联 $X$ 与 $Y$ 的观测时，前者是对称的（例如观测到 $Y=0$ 蕴含 $\beta x=-\varepsilon$），而在涉及干预时，它们则是不对称的（例如把 $Y$ 设为零，并不能告诉我们 $x$ 与 $\varepsilon$ 之间的关系）。路径图中的箭头把这种双重角色明确表示出来，这也许正是使用图所获得的洞见与推断力的来源。

方程 $y=\beta x+\varepsilon$ 最强的经验断言，是通过把其他变量排斥在方程右端之外而作出的，从而宣称 $X$ 是 $Y$ 的\emph{唯一}直接原因。这转化为一个可检验的不变性断言：在条件 $\dox{x}$ 下 $Y$ 的统计量，应当对模型中任何其他变量的操纵保持不变（见 \S 1.3.2）。\footnote{结构方程对系统某些变化保持不变这一基本观念，可上溯至 Marschak（1950）与 Simon（1953），并在 Hurwicz（1962）、Mesarovic（1969）、Sims（1977）、Cartwright（1989）、Hoover（1990）与 Woodward（1995）中以各种抽象层次获得了数学表述。然而，\textup{(5.23)} 的简明、精确与清晰是无与伦比的。}这一断言可以符号地写作
\begin{equation}
P(y\mid \dox{x},\dox{z})=P(y\mid \dox{x}) \tag{5.23}
\end{equation}
对一切与 $\{X\cup Y\}$ 不相交的 $Z$ 成立。\footnote{事实上，这一断言只是方程所传达之信息的一部分；另一部分是动态的或反事实的断言：若我们把 $X$ 控制到 $x'$ 而非 $x$，则 $Y$ 将取值 $\beta x'+\varepsilon$。换言之，在 $X$ 的各种假想控制下、并在相同的外部条件 $(\varepsilon)$ 下描点 $Y$ 的取值，应当得到一条斜率为 $\beta$ 的直线。这类关于系统在相继控制条件下之行为的确定性动态断言，只有在如下假设下才可检验：$\varepsilon$——它代表外部条件或实验单位的属性——在我们从 $x$ 切换到 $x'$ 时保持不变。这类反事实断言构成了每一条科学定律的经验内容（见 \S 7.2.2）。}相反，回归方程不作任何经验断言。

注意，这一不变性是相对于对 $Z$ 的\emph{操纵}、而非对 $Z$ 的\emph{观测}而言的。在条件 $\dox{x}$ 下并给定测量 $Z=z$ 时 $Y$ 的统计量，记作 $P(y\mid \dox{x},z)$；若该测量取自 $Y$ 的某个后果（即后代），它当然会依赖于 $z$。还要注意，普通条件概率 $P(y\mid x)$ 并不享有如此强的不变性，因为 $P(y\mid x)$ 一般会对模型中 $X$ 以外变量的操纵敏感（除非 $X$ 与 $\varepsilon$ 独立）。相反，式 \textup{(5.23)} 不论 $\varepsilon$ 与 $X$ 之间的统计关系如何都成立。

推广到由若干结构方程组成的方程组，\textup{(5.23)} 阐明了给定因果图背后的假设。若 $G$ 是与一组结构方程相联系的图，则这些假设以下述方式体现于 $G$ 中：(1) 每一条缺失的箭头——比如 $X$ 与 $Y$ 之间者——代表如下假设：一旦我们实施干预并固定 $Y$ 的父节点，$X$ 对 $Y$ 便无因果效应；(2) 每一条 $X$ 与 $Y$ 之间缺失的双向连线，代表如下假设：那些（直接）影响 $X$ 的遗漏因子，与那些（直接）影响 $Y$ 的遗漏因子不相关。后一假设的操作含义，我们将在 \textup{(5.25)}—\textup{(5.27)} 中定义。

\subsubsection*{结构参数：操作定义}

把结构方程解释为关于 $Y$ 在假想干预下之行为的陈述，为结构参数提供了一个简单的定义。方程 $y=\beta x+\varepsilon$ 中 $\beta$ 的含义，就是
\begin{equation}
\beta=\frac{\partial}{\partial x}E[Y\mid \dox{x}], \tag{5.24}
\end{equation}
即在 $X$ 由外部控制保持在 $x$ 的实验中，$Y$ 的期望（相对于 $x$）的变化率。这一解释不论在非实验研究中 $\varepsilon$ 与 $X$ 是否相关（例如经由另一方程 $x=\alpha y+\delta$）都成立。

至此我们几乎无需补充：$\beta$ 与回归系数 $r_{YX}$ 毫无关系，等价地说，与条件期望 $E(Y\mid x)$ 毫无关系——尽管许多教科书如此暗示。$\beta$ 与回归系数相重合的条件，已在定理 5.3.1 中列明。

尽管如此，把 \textup{(5.24)} 的定义与那些承认 $\beta$ 之不变特征、却难以说清 $\beta$ 对何种变化保持不变的理论相比较，仍很重要。例如，Cartwright（1989, p.~194）把 $\beta$ 刻画为一种她称为“能力”（capacity）的自然不变量。她正确地指出，$\beta$ 在变化中保持不变，但接着解释说：随着 $X$ 的统计量改变，“保持不变的是比值 $[\beta=E(YX)/E(X^2)]$，无论方差如何移动”。这一刻画在两方面不精确。第一，$\beta$ 一般可能既不等于所说的比值，也不等于统计参数的任何其他组合。第二——这也是定义 5.4.1 的要点——结构参数对\emph{局部}干预（即系统内特定方程的改变）不变，而对变量统计量的一般变化并不不变。若我们从 $\mathrm{cov}(X,\varepsilon)=0$ 出发，而 $X$ 的方差之所以改变，是因为我们（或自然）局部地修改了生成 $X$ 的过程，那么 Cartwright 是对的：比值 $\beta=E(YX)/E(X^2)$ 将保持不变。然而，若 $X$ 的方差因任何其他原因而改变——比如说，因为我们观测到某个同时依赖 $X$ 与 $Y$ 的证据 $Z=z$，或者因为生成 $X$ 的过程变得依赖于更广的一组变量——则该比值不会保持不变。

\subsubsection*{神秘的误差项：操作定义}

定义 5.4.1 与 \textup{(5.24)} 所给的解释，为那个神秘的误差项提供了操作定义：
\begin{equation}
\varepsilon=y-E[Y\mid \dox{x}], \tag{5.25}
\end{equation}
它在非操纵性研究中虽不可观测，却远非某些研究者所暗示的那样是形而上的或纯定义的（如 Richard 1980; Holland 1988, p.~460; Hendry 1995, p.~62）。与回归方程中的误差不同，$\varepsilon$ 度量的是 $Y$ 对其\emph{受控}期望 $E[Y\mid \dox{x}]$ 的偏离，而不是对其条件期望 $E[Y\mid x]$ 的偏离。因此，一旦 $X$ 受到控制，$\varepsilon$ 的统计量便可由对 $Y$ 的观测测得。或者，由于不论 $X$ 被操纵还是被观测 $\beta$ 都保持不变，若已知 $\beta$，$\varepsilon=y-\beta x$ 的统计量便可在观测研究中测得。

同样，误差之间的相关也可以经验地估计。对任意两个不相邻的变量 $X$ 与 $Y$，\textup{(5.25)} 给出
\begin{equation}
E[\varepsilon_Y\varepsilon_X]=E[YX\mid \dox{\pa_Y,\pa_X}]-E[Y\mid \dox{\pa_Y}]E[X\mid \dox{\pa_X}]. \tag{5.26}
\end{equation}
一旦我们确定了结构系数，受控期望 $E[Y\mid \dox{\pa_Y}]$、$E[X\mid \dox{\pa_X}]$ 与 $E[YX\mid \dox{\pa_Y,\pa_X}]$ 便成为观测变量 $\pa_Y$ 与 $\pa_X$ 的已知线性函数；因此，\textup{(5.26)} 右端的期望可以在观测研究中估计。另一方面，若系数尚未确定，则该表达式可以在干预研究中直接评估：固定 $\pa_X$ 与 $\pa_Y$（假定 $X$ 与 $Y$ 不处于父子关系），并由如此条件下所得的数据估计 $X$ 与 $Y$ 的协方差。

最后，我们常常并不关心评估 $E[\varepsilon_Y\varepsilon_X]$ 的数值，而关心判定是否可以假定 $\varepsilon_Y$ 与 $\varepsilon_X$ 不相关。为此判定，只需检验等式
\begin{equation}
E[Y\mid x,\dox{s_{XY}}]=E[Y\mid \dox{x},\dox{s_{XY}}] \tag{5.27}
\end{equation}
是否成立；其中 $s_{XY}$ 代表模型中除 $X$ 与 $Y$ 之外一切变量的（任一）设定。这一检验可以施于模型中任意两个变量，唯有当 $Y$ 是 $X$ 的父节点时例外——此时适用将 $X$ 与 $Y$ 互换的对称方程。

\subsubsection*{神秘的误差项：概念解释}

SEM 教科书的作者通常把误差项解释为遗漏因子的影响。许多 SEM 研究者不愿接受这一解释，部分是因为未指明的遗漏因子为形而上的思辨打开了大门，部分是因为基于这类因子的论证曾被不当利用，成为从路径图中省略双向弧的通用而空洞的许可（McDonald 1997）。这些忧虑可由误差项的操作解释 \textup{(5.25)} 来回答，因为它规定的是误差如何被\emph{测量}，而不是它们如何\emph{起源}。

然而必须注意：在判定是否可以假定成对的误差项不相关时，这一操作定义并不能取代遗漏因子的构想。因为这类判定需要在模型参数仍然“自由”的阶段作出，它们不可能建立在相关的数值评估之上，而必须立足于定性的结构知识——关于机制如何彼此勾连、变量如何相互影响的知识。这种判断性的决定，几乎无法由 \textup{(5.26)} 的操作判据来辅助——该判据要求研究者评估：在复杂的实验条件下，就两个不同变量所取的两个偏离，是否会有相关。这样的评估在认知上是不可行的。

相反，遗漏因子的构想要求研究者判断：是否存在某些因子同时影响若干观测变量。这样的判断在认知上是可以处理的，因为它们是定性的，并且立足于纯然结构性的知识——而这正是在建模的这一阶段唯一可用的知识。

研究者还应考虑误差相关的另一来源：选择偏倚。若两个不相关的未观测因子具有一个共同效应，该效应被排除在分析之外、却影响了本研究的样本选择，则在被抽样的总体中相应的误差项将相关；因此，\textup{(5.26)} 中的期望在被抽样总体上取时不会消失（参见 \S 1.2.3 关于伯克森悖论的讨论）。

然而我们应当强调：图中\emph{缺失}的弧，而非图中已有的弧，最需要关注与审慎的实质论证。多加一条双向弧，最坏不过是损害参数的可识别性；而删去一条已有的双向弧，则可能产生错误的结论，并带来一种虚假的模型可检验感。因此，应当默认图中任意两个节点之间存在双向弧；只有具备充分理据时——例如这两个变量不太可能存在共同原因、也不太可能存在选择偏倚——才应删去它们。尽管我们永远无法知晓可能影响我们变量的全部因子，实质知识有时允许我们断言：某个可能的共同因子的影响不大可能显著。

于是，正如科学中常有的情形，我们测量物理实体的方式，并不提供思考它们的最佳方式。误差的遗漏因子构想立足于结构知识，因而在构建、评估与思考因果模型时，是比操作定义更有用的指南。

\subsection{效应分解的解释}

结构方程建模引以为荣（而且理当如此）的是：它提供了区分直接效应与间接效应的有原则的方法论。我们在 \S 4.5 中已经看到，这种区分在从过程控制到法律纠纷的许多应用中都很重要，而 SEM 确实提供了定义、识别与估计直接效应和间接效应的一套自洽方法论。然而，多数 SEM 研究者不愿承认结构参数的因果解读——加之他们热衷于代数运算——导致了直接效应与间接效应定义不当，正如 Freedman（1987）与 Sobel（1990）所指出的。本节希望借助对结构系数之操作含义的恪守，来纠正这一混乱。

我们从定义 3.2.1 中因果效应 $P(y\mid \dox{x})$ 的一般概念出发，然后把它特殊化以定义直接效应（如 \S 4.5 那样），最后把这些定义用结构系数表达出来。

\begin{definition}[总效应]
$X$ 对 $Y$ 的\emph{总效应}由 $P(y\mid \dox{x})$ 给出，即：当 $X$ 保持在 $x$、而其他一切变量任其自然运行时 $Y$ 的分布。
\end{definition}

\begin{definition}[直接效应]
$X$ 对 $Y$ 的\emph{直接效应}由 $P(y\mid \dox{x},\dox{s_{XY}})$ 给出，其中 $S_{XY}$ 是系统中除 $X$ 与 $Y$ 之外一切观测变量之集合。
\end{definition}

在线性分析中，定义 5.4.2 与定义 5.4.3 对 $x$ 求导之后，给出通常用以定义直接效应与间接效应的那些熟悉的路径系数。然而，它们在若干重要方面不同于常规定义。第一，直接效应是用假想实验来定义的：在那些实验中，中间变量由物理干预保持不变，而不是由统计调整（后者常被藏在“控制住”这一误导性说法之下）保持不变。图 5.10 给出一个简单的例子：对中间变量（$Z$ 与 $W$）作调整不会给出 $X$ 对 $Y$ 之直接效应的正确值零，而 $\partial E(Y\mid \dox{x,z,w})/\partial x$ 却给出正确的值：
\[
\frac{\partial}{\partial x}(\alpha w+\gamma z)=0.
\]
\S 4.5.3（表 4.1）提供了另一个这样的例子，其中涉及二分变量。

第二，没有必要把控制仅限于中间变量；系统中一切变量都可以保持不变（$X$ 与 $Y$ 除外）。在假想的意义上，科学家控制一切可能的条件 $S_{XY}$，而测量可以在尚不知图的结构的情况下开始。最后，我们的定义与常规不同之处还在于：它把总效应与直接效应彼此独立地解释，作为两个不同实验的结果。教科书定义（如 Bollen 1989, p.~376; Mueller 1996, p.~141; Kline 1998, p.~175）通常把总效应等同于路径系数矩阵的幂级数。这一代数定义在递归（半马尔可夫）系统中与操作定义（定义 5.4.2）一致，但在含反馈的模型中则给出错误的表达式。例如，给定方程组 $\{y=\beta x+\varepsilon,\ x=\gamma y+\delta\}$，$X$ 对 $Y$ 的总效应就是 $\beta$，而不是 Bollen（1989, p.~379）所断言的 $\beta(1-\gamma\beta)^{-1}$。后者并不具有配得上“$X$ 的效应”这一说法的操作意义。\footnote{这一错误由 Sobel（1990）指出；但也许因为路径系数之恒定性被当作一个新的、外来的假设提出，Sobel 的更正并未带来实践或哲学上的转变。}

在结束本节关于效应分解的讨论之前，我们还有几句关于二分变量的话，应当会引起处理此类变量的研究者的兴趣。这类变量之间的关系通常是非线性的，因此 \S 4.5 的结果应当适用。特别地，$X$ 对 $Y$ 的直接效应将依赖于我们把 $Y$ 的其他父节点保持在什么水平。若我们希望对这些取值取平均，便得到 \textup{(4.11)} 所给的表达式——它诉诸嵌套反事实，并可化简为 \textup{(4.12)}。

我们在定义结构方程之经验内容（定义 5.4.1）时所诉诸的操纵论说明，在线性系统中是足够的——那里多数所关心的因果量都可以由总体层面的实验研究推断（见 \S 11.7.1）。在非线性与非参数模型中，我们有时需要下降到个体单位层面，诉诸结构方程更基本的反事实解读，如式 \textup{(3.51)} 与第 160 页注 17 所阐述。间接效应的分析即是一例；其定义 \textup{(4.14)} 立足于嵌套反事实，无法用总体平均来表达。为使间接效应获得独立于总效应与直接效应的操作含义，这种分析是必需的（见 \S 11.4.2）。不过，借助反事实语言，我们可以给间接效应一个简单的操作定义：$X$ 对 $Y$ 的间接效应，就是在保持 $X$ 不变、并把中介变量 $Z$ 增至“$X$ 增加一个单位时 $Z$ 本会达到的取值”时，我们将会看到的 $Y$ 的增量（形式定义见 \S 4.5.5）。在线性系统中，这一定义确实与总效应和直接效应之差相合。关于间接效应在社会科学与政策分析中的作用，进一步讨论见第十一章（Pearl 2005a）。

\subsection{外生性、超外生性及其他枝节}

经济学教科书总是不忘提醒读者：外生变量与内生变量之间的区分，一方面“对建模最为重要”（Darnell 1994, p.~127），另一方面又是“一个微妙而时有争议的复杂问题”（Greene 1997, p.~712）。经济学学生自然会期待本章所发展的概念与工具能对这一问题有所启迪——他们理当如此期待。下面我们给出一个简单的外生性定义，它抓住了文献中出现的那些重要细微差别，既易懂又精确。

如今流行的做法是区分三种外生性：弱的、强的与超的（Engle et al. 1983）；前两者是统计性的，后者是因果性的。然而，外生性之所以重要——也之所以地位有争议——在于它对政策干预的蕴含。因此有些经济学家认为，只有因果方面（即超外生性）才配得上“外生的”这一名号，而统计性的版本则是不正当的闯入者，往往把识别性与可解释性的问题同估计效率的问题混为一谈（Ed Leamer，私人通信）。\footnote{John Aldrich 与 James Heckman 也表达了类似的意见。另见 Aldrich（1993）。}为同时服务两个阵营，我将先给出因果外生性的简单定义，然后给出一个更一般的定义，使因果方面与统计方面都作为特例从中导出。于是，我们所称的“外生性”对应于 Engle 等所称的“超外生性”，这一概念把握了某些关系在政策干预下的结构不变性。

假定我们考虑对一组变量 $X$ 实施干预，并希望刻画一组结局变量 $Y$ 在干预 $\dox{X=x}$ 下的统计行为。用通常的记号 $P(y\mid \dox{x})$ 表示 $Y$ 的干预后分布。若我们感兴趣的是该分布的某一组参数 $\lambda$，则任务是从可得数据中估计 $\lambda[P(y\mid \dox{x})]$。然而，可得数据通常是在另一组条件下生成的：$X$ 并未被保持不变，而是随那些促使决策者过去设定 $X$ 的经济压力与预期而变动。用 $M$ 表示过去生成数据的过程，用 $P_M(y)$ 表示与 $M$ 相联系的概率分布，我们便问：在关于 $M$ 的背景知识 $T$（意为“理论”）之下，$\lambda[P_M(y\mid \dox{x})]$ 能否由从 $P_M(y)$ 抽取的样本一致地估计？这实质上就是我们本章及前几章所分析的识别问题，但有一个重要区别：我们现在问的是，$\lambda[P(y\mid \dox{x})]$ 能否仅由\emph{条件}分布 $P(y\mid x)$ 识别，而不是由整个联合分布 $P(y)$ 识别。当识别在这一受限条件下成立时，便称 $X$ 相对于 $(Y,\lambda,T)$ 是\emph{外生的}。

我们可以把这形式地陈述如下。

\begin{definition}[外生性]
设 $X$ 与 $Y$ 是两个变量集，$\lambda$ 是干预后概率 $P(y\mid \dox{x})$ 的任一参数集。若 $\lambda$ 可由条件分布 $P(y\mid x)$ 识别，即若
\begin{equation}
P_{M_1}(y\mid x)=P_{M_2}(y\mid x)\ \Longrightarrow\ \lambda[P_{M_1}(y\mid \dox{x})]=\lambda[P_{M_2}(y\mid \dox{x})] \tag{5.28}
\end{equation}
对任何两个满足理论 $T$ 的模型 $M_1$ 与 $M_2$ 成立，则称 $X$ 相对于 $(Y,\lambda,T)$ 是外生的。
\end{definition}

在 $\lambda$ 构成干预后概率之完整刻画的特殊情形中，\textup{(5.28)} 化简为蕴含
\begin{equation}
P_{M_1}(y\mid x)=P_{M_2}(y\mid x)\ \Longrightarrow\ P_{M_1}(y\mid \dox{x})=P_{M_2}(y\mid \dox{x}). \tag{5.29}
\end{equation}
若我们进一步假定：对每个 $P(y\mid x)$，我们的理论 $T$ 并不先验地排除某个满足 $P_{M_2}(y\mid \dox{x})=P_{M_2}(y\mid x)$ 的模型 $M_2$，\footnote{例如，若 $T$ 代表一切具有相同图结构的模型，则这样的 $M_2$ 并不被先验排除。}则 \textup{(5.29)} 化简为等式
\begin{equation}
P(y\mid \dox{x})=P(y\mid x), \tag{5.30}
\end{equation}
我们认出这一条件就是“无混杂”（见 \S 3.3 与 \S 6.2）。式 \textup{(5.30)} 由 \textup{(5.29)} 推出，因为 \textup{(5.29)} 必须对 $T$ 中所有 $M_1$ 成立。注意，由于理论 $T$ 未被显式提及，\textup{(5.30)} 可以施于任何单个模型 $M$，并可当作外生性的又一个定义——尽管它比 \textup{(5.28)} 更强。

坚持要求 $\lambda$ 仅由条件分布 $P(y\mid x)$ 识别（尽管边缘分布 $P(x)$ 也可得），其动机在于这对估计过程的蕴含。如 \textup{(5.30)} 所述，发现 $X$ 是外生的，使我们能够直接由被动观测预测（对 $X$ 的）干预效果，甚至无需对混杂因子作调整。我们在 \S 3.3 与 \S 5.3 中的分析进一步提供了外生性的图形检验：若不存在从 $X$ 到 $Y$ 的未阻断后门路径，则 $X$ 关于 $Y$ 是外生的（定理 5.3.2）。这一检验以程序性定义补充了 \textup{(5.30)} 的陈述性定义，从而完成了外生性的形式化。通常被归诸超外生性的那些不变性可以从因果图的拓扑中辨认出来，这不应令人惊讶——考虑到每张因果图都表示一个结构模型，而每个结构模型都已体现了政策预测所需的不变性假设（见定义 5.4.1）。

Leamer（1985）把 $X$ 定义为外生的，若 $P(y\mid x)$ 对“生成”$X$ 之过程的变化保持不变。这一定义与 \textup{(5.30)} 相合，\footnote{前提是变化仅限于函数的修改，而不改变各函数中自变量的集合（即父节点）。}因为 $P(y\mid \dox{x})$ 由这样一个结构模型支配：其中决定 $X$ 的方程被抹去；因此 $P(y\mid x)$ 必定对那些方程的性质不敏感。相反，Engle 等（1983）用 $X$ 的“边缘密度”的变化来定义外生性（即他们的超外生性）；一如往常，从过程语言转向统计术语会导致歧义。据 Engle 等（1983, p.~284），外生性要求条件分布 $P(y\mid x)$ 的一切参数“对条件变量分布的任何变化都不变”，\footnote{这一要求在 Darnell（1994, p.~131）与 Maddala（1992, p.~192）中被逐字重复。}（即 $P(x)$ 的任何变化）。这种“对 $P(x)$ 的任何变化都保持不变”的要求过强——即便 $X$ 完全外生，变化的条件或新的观测也很容易同时改变 $P(x)$ 与 $P(y\mid x)$。（举例来说，考虑一种把 $X$ 无可争议地外生的随机化实验变成非随机化实验的变化；我们不应当坚持 $P(y\mid x)$ 在这种变化下保持不变。）所考虑的变化类必须被限制为对决定 $X$ 之机制（或方程）的局部修改，如 Leamer 所述；而这一限制必须被纳入任何外生性定义之中。然而，为使这一限制精确，必须像 $P(y\mid \dox{x})$ 的定义那样诉诸 SEM 的词汇；边缘密度与条件密度的词汇太过粗糙，不足以恰当地定义 $P(y\mid x)$ 应当对之保持不变的那些变化。

现在我们准备定义一个更一般的外生性概念，把“弱”外生性与“超”外生性纳入同一框架。\footnote{我们略去对“强”外生性的讨论，它是弱外生性在时间序列分析中适用的一个稍为繁复的版本。}为此，我们从定义 5.4.4 中除去“$\lambda$ 必须表示干预后分布之特征”这一限制。相反，我们允许 $\lambda$ 表示底层模型 $M$ 的任何特征，包括结构特征（如路径系数、因果效应与反事实），也包括统计特征（当然，后者可以仅由联合分布确定）。有了这一推广，我们同时也得到了一个更简单的外生性定义。

\begin{definition}[一般外生性]
设 $X$ 与 $Y$ 是两个变量集，$\lambda$ 是理论 $T$ 中某个结构模型 $M$ 上定义的任一参数集。若 $\lambda$ 可由条件分布 $P(y\mid x)$ 识别，即若
\begin{equation}
P_{M_1}(y\mid x)=P_{M_2}(y\mid x)\ \Longrightarrow\ \lambda(M_1)=\lambda(M_2) \tag{5.31}
\end{equation}
对任何两个满足理论 $T$ 的模型 $M_1$ 与 $M_2$ 成立，则称 $X$ 相对于 $(Y,\lambda,T)$ 是外生的。
\end{definition}

当 $\lambda$ 由结构参数（如路径系数或因果效应）组成时，\textup{(5.31)} 表达的是对种种干预（而不仅是 $\dox{X=x}$）的不变性。尽管干预本身未在 \textup{(5.31)} 中被显式提及，等式 $\lambda(M_1)=\lambda(M_2)$ 却通过 $\lambda$ 的结构性质反映了这类干预。特别地，若 $\lambda$ 代表因果效应函数 $P(y\mid \dox{x})$ 在 $x$ 与 $y$ 的选定点上的取值，则 \textup{(5.31)} 化简为蕴含
\begin{equation}
P_{M_1}(y\mid x)=P_{M_2}(y\mid x)\ \Longrightarrow\ P_{M_1}(y\mid \dox{x})=P_{M_2}(y\mid \dox{x}), \tag{5.32}
\end{equation}
它与 \textup{(5.29)} 完全相同。于是外生性的因果性质随之而得。

当 $\lambda$ 由严格统计性的参数——如均值、众数、回归系数或其他分布特征——组成时，$M$ 的结构特征不进入考虑；我们有 $\lambda(M)=\lambda(P_M)$，于是 \textup{(5.31)} 化简为
\begin{equation}
P_1(y\mid x)=P_2(y\mid x)\ \Longrightarrow\ \lambda(P_1)=\lambda(P_2) \tag{5.33}
\end{equation}
对任何两个与 $T$ 相容的概率分布 $P_1(x,y)$ 与 $P_2(x,y)$ 成立。我们由此得到了一个统计性的外生性概念，它允许我们在估计 $\lambda$ 时忽略边缘分布 $P(x)$；我们可以称之为“弱外生性”。\footnote{Engle 等（1983）还施加了一个称为“变差自由”（variation-free）的要求；在处理真正的结构模型 $M$（其中各机制之间互不约束）时，这一要求自动满足。}

最后，若 $\lambda$ 由 $Y$ 中变量之间（不含 $X$）的因果效应组成，我们便得到工具变量的一个推广定义。例如，若我们关心的是因果效应 $\lambda=P(w\mid \dox{z})$，其中 $W$ 与 $Z$ 是 $Y$ 中的两个变量集，则 $X$ 相对于该参数的外生性保证了 $P(w\mid \dox{z})$ 可由条件概率 $P(z,w\mid x)$ 识别。这正是工具变量的作用——协助识别不涉及该工具的因果效应。（参见图 5.9，其中 $Z,X,Y$ 分别代表 $X,Z,W$。）

关于多数教科书所用语言，还有一句告诫：外生性常被定义为“参数是否‘进入’条件密度或边缘密度的表达式”。例如，Maddala（1992, p.~392）把弱外生性定义为：边缘分布 $P(x)$“不涉及”$\lambda$。这类定义并不无歧义，因为“参数是否进入某个密度”或“某个密度是否涉及某个参数”这类问题依赖于句法；不同的代数表示可能使某些参数显明，也可能使之隐晦。例如，若 $X$ 与 $Y$ 是二分变量，所关心的参数是 $\lambda=P(y_0\mid x_0)$，则边缘概率 $P(x)$ 当然“涉及”诸如
\[
\lambda_1=P(x_0,y_0)+P(x_0,y_1)
\quad\text{与}\quad
\lambda_2=P(x_0,y_0)
\]
这样的参数，而 $\lambda$“涉及”它们的比值：
\[
\lambda=\lambda_2/\lambda_1.
\]
因此，写出 $P(x_0)=\lambda_2/\lambda$，便表明 $\lambda$ 与 $\lambda_2$ 都涉入边缘概率 $P(x_0)$，人们也许会由此断定 $X$ 相对于 $\lambda$ 不是外生的。然而，$X$ 事实上相对于 $\lambda$ 是外生的，因为比值 $\lambda=\lambda_2/\lambda_1$ 不是别的，正是 $P(y_0\mid x_0)$；因此它如 \textup{(5.33)} 所要求的那样，由 $P(y_0\mid x_0)$ 唯一确定。\footnote{Engle 等（1983, p.~281）与 Hendry（1995, pp.~162—163）通过有选择的“重参数化”克服了这一歧义——这是一个必要步骤，而教科书往往忽略它。}

\textup{(5.31)} 所给定义的优点在于：它不依赖于密度函数的句法表示，而仅依赖其语义内容。参数被当作从模型计算出来的量，而不是描述模型的数学符号。因此，这一定义既适用于统计参数，也适用于结构参数；事实上，它适用于任何可以从结构模型 $M$ 计算出来的量 $\lambda$，不论该量是否（或可能）用于描述边缘密度或条件密度。

\subsubsection*{再论神秘的误差项}

从历史上看，引起争议最多的外生性定义，是用变量与误差之间的相关来表达的那一个。它读作：

\begin{definition}[基于误差的外生性]
变量 $X$ 是外生的（相对于 $\lambda=P(y\mid \dox{x})$），若 $X$ 独立于一切影响 $Y$ 的误差，但以 $X$ 为中介者除外。
\end{definition}

这一定义由 Hendry and Morgan（1995）上溯至 Orcutt（1952），在 1950 至 1970 年间成为计量经济学文献的标准（如 Christ 1966, p.~156; Dhrymes 1970, p.~169），至今仍引导着多数计量经济学家的思考（例如在工具变量的选择中；Bowden and Turkington 1984）。然而，它在 20 世纪 80 年代初遭到批评，当时结构误差（式 \textup{(5.25)}）与回归误差之间的区分变得模糊（Richard 1980）。（回归误差按定义与回归变量正交。）Cowles 委员会关于结构方程的逻辑（见 \S 5.1）尚未达到完全的数学成熟；由于否认结构参数与回归参数之间的记号区分，它使一切基于误差项的表述都蒙上了歧义的嫌疑。建立全新外生性基础的前景——表面上摆脱了“误差”与“结构”这类理论术语（Engle et al. 1983）——进一步使经济学家不愿去整理 Cowles 委员会的逻辑，而对基于误差之外生性定义的批评则日益时髦。例如，Hendry and Morgan（1995）写道：“外生性概念迅速蜕化为一个松散的观念：一个可观测变量具有与某个未观测误差不相关的性质”；而 Imbens（1997）欣然同意这一观念“是不充分的”。\footnote{Imbens 更偏爱以“随机分配假设”这类实验隐喻来下定义，也许是因为担心“研究者通常对这些人扰动究竟代表什么并无确定的想法”（Angrist et al. 1996, p.~446）。我不赞同；“随机分配”是一个误导性的隐喻，而“遗漏因子”则清晰透彻。}

如果我们考虑在使用恰当记号（如 \textup{(5.25)}）时结构误差可以被定义得何等精确清晰，这些批评者便谈不上有道理。当施于结构误差时，标准的基于误差的外生性判据在形式上与 \textup{(5.30)} 相合，这可以用定理 5.3.2 的后门检验来验证（取 $Z=0$）。因此，标准定义所传达的信息，与那些更复杂、更难沟通的外生性定义所蕴含的信息相同。因此我确信：标准定义终将重新赢回它一直理应享有的接受与体面。

外生性与工具变量之图形定义与反事实定义之间的关系，将在第七章（\S 7.4.5）讨论。

\section{结论}

今天，被称为结构方程建模的这一事业正日益受到攻击。创立者们已经退休，他们的教导已被遗忘，而实践者、教师与研究者发现，他们继承下来的这套方法论既难以辩护，也难以取代。现代 SEM 教科书专注于参数估计，很少阐明这些参数在因果解释或政策分析中所扮演的角色；例如，讨论干预效果的例子显著缺席。SEM 的研究如今几乎完全集中于模型拟合，而涉及 SEM 模型之意义与用法的问题，则成为混乱与争议的题材。其中一些混乱反映在我收到的读者提问中（\S 11.5）；为此我为他们奉献了一份“SEM 生存包”（\S 11.5.3）——一组为 SEM 的因果解读及其科学理据作辩护的论证。

我完全确信，SEM 的当代危机源于缺乏一套处理结构方程中因果信息的数学语言。图模型提供了这样一套语言。它们因而帮助我们回答了驱动当前危机的许多悬而未决的问题：

\begin{enumerate}[label=\arabic*.,leftmargin=1.8em]
\item 在什么条件下，我们可以赋予结构系数以因果解释？
\item 某个给定结构方程模型背后的因果假设是什么？
\item 任何给定结构方程模型的统计蕴含是什么？
\item 某个给定结构系数的操作含义是什么？
\item 任何给定结构方程模型作出了哪些政策制定方面的断言？
\item 一个方程何时不是结构的？
\end{enumerate}

本章描述了如今解决这些基础性问题的概念进展。（\S 11.5.2 与 \S 11.5.3 有进一步阐述。）此外，我们提出了若干工具，用以回答具有实践重要性的问题：

\begin{enumerate}[label=\arabic*.,leftmargin=1.8em]
\item 两个结构方程模型何时在观测上不可区分？
\item 回归系数何时代表路径系数？
\item 加入一个回归变量何时会引入偏倚？
\item 在收集任何数据之前，我们如何判断哪些路径系数可以被识别？
\item 何时可以弃置线性—正态假设，却仍从数据中提取因果信息？
\end{enumerate}

我仍抱有希望：研究者会认识到这些概念与工具的益处，并借助它们重振社会与行为科学中的因果分析。

\section{第二版后记}

\subsection{计量经济学的觉醒？}

在经济学中因果分析被忽视数十年之后，一股兴趣的涌动似乎正在发生。在最近一系列论文中，Jim Heckman（2000, 2003, 2005, 2007（与 Vytlacil））作出了巨大努力，以复活并重申 Cowles 委员会对结构方程模型的解释，并使经济学家相信：因果分析的最新进展根源于 Haavelmo（1943）、Marschak（1950）、Roy（1951）与 Hurwicz（1962）的思想。遗憾的是，Heckman 仍未向计量经济学家提供本章所提问题（第 133、170、171、215—217 页）的清晰答案。特别地，Heckman 过分关注实施问题，拒绝把 Haavelmo 的“方程抹除”作为定义反事实的基础，却未能向计量经济学家提供另一种定义——即像式 \textup{(3.51)} 那样，在一个适定的经济模型中计算反事实 $Y(x,u)$ 的程序，其中 $X$ 与 $Y$ 是模型中任意两个变量。（见 \S 11.5.4—5。）这样的定义对于赋予“潜在结局”进路以基于 SEM 的形式语义、从而统一目前彼此隔绝的两个计量经济学阵营，是不可或缺的。

积极觉醒的另一个迹象来自社会科学，即 Morgan 与 Winship 的著作《反事实与因果推断》（Counterfactual and Causal Inference, 2007）的出版，其中 SEM 的因果解读被明确恢复。\footnote{不过，反事实的 SEM 基础遗憾的是未获阐明。}

\subsection{线性模型中的识别}

在一系列论文中，Brito and Pearl（2002a,b, 2006）建立了若干图形判据，把可识别的半马尔可夫线性模型类显著扩展到超出本章所讨论的范围。他们首先证明：在一切不含弓形弧的图中——即不允许一个原因与其直接效应之间存在误差相关，而对与间接原因相联系的误差不施加任何限制——识别都得到保证（Brito and Pearl 2002b）。随后，他们把工具变量的概念推广到图 5.9 与图 5.11 的经典模式之外，建立了一个可在多项式时间内检验、且囊括文献中一切已知条件的通用识别条件。代数进路见 McDonald（2002a），浅近介绍与结果综述见 Brito（2010）。

\subsection{因果断言的稳健性}

SEM 中的因果断言，是由数据与模型所体现的一组因果假设共同确立的。例如，“图 5.9 中因果效应 $E(Y\mid \dox{x})$ 由 $\alpha x=r_{YZ}/r_{XZ}\,x$ 给出”这一断言，基于如下假设：$\mathrm{cov}(\varepsilon_Z,\varepsilon_Y)=0$ 与 $E(Y\mid \dox{x,z})=E(Y\mid \dox{x})$；二者都在图中显示出来。一个断言是\emph{稳健的}，当它对模型中某些假设的违反不敏感。例如，上述断言对假设 $\mathrm{cov}(\varepsilon_Z,\varepsilon_X)=0$ 不敏感，而该假设在模型中也显示出来。

当若干组不同的假设给出参数 $\alpha$ 的 $k$ 个不同估计量时，称该参数是 $k$-\emph{可识别的}；$k$ 越大，基于 $\alpha$ 的断言就越稳健，因为这些估计量之间的相等对协方差矩阵施加了 $k-1$ 个约束；若这些约束在数据中得到满足，便表明 $k$ 组不同假设彼此一致，从而支持它们的有效性。一个典型的例子出现在有若干（独立的）工具变量 $Z_1,Z_2,\dots,Z_k$ 可用于单一连线 $X\to Y$ 之时，它们给出等式
\[
\alpha=r_{YZ_1}/r_{XZ_1}=r_{YZ_2}/r_{XZ_2}=\cdots=r_{YZ_k}/r_{XZ_k}.
\]
Pearl（2004）给出了这一稳健性概念的形式定义，并建立了量化给定因果断言之稳健程度的图形条件。$k$-识别推广了标准 SEM 分析中自由度的概念；后者刻画整个模型，而前者适用于个别参数，更一般地说，适用于个别的因果断言。

\section*{致谢}
\addcontentsline{toc}{section}{致谢}

本章的灵感，得自于一代又一代统计学家——他们带着幽默与怀疑追问：SEM 的方法论怎么可能对任何理性存在者有意义；也得自于那些社会科学家——他们（也许是不自觉地）把 SEM 传统从统计学解释的淹没中拯救了出来。Herman Ader、Peter Bentler、Kenneth Bollen、Jacques Hagenaars、Rod McDonald、Les Hayduk 与 Stan Mulaik 的评论，帮助我更深入地理解 SEM 的实践与词汇。John Aldrich、Nancy Cartwright、Arthur Goldberger、James Heckman、Kevin Hoover、Ed Leamer 与 Herbert Simon 帮助我穿透了计量经济学中结构方程与外生性的迷宫。金天（Jin Tian）在修订 \S 5.2.3 与 \S 5.3.1 时贡献良多。
